\documentclass[11pt]{article}

\usepackage{fancyhdr}
\usepackage{style}
\usepackage{enumitem}
\usepackage{hyperref}
\usepackage[yyyymmdd]{datetime}
\usepackage[letterpaper, margin=0.5in, top=0.8in]{geometry}

\newdateformat{mydate}{\shortmonthname[\THEMONTH] \THEYEAR}
\mydate

\newcommand{\daterange}[4]{\formatdate{1}{#1}{#2} - \formatdate{1}{#3}{#4}}
\newcommand{\daterangepres}[2]{\formatdate{1}{#1}{#2} - Present}

\setlength\parindent{0pt}

\pagenumbering{gobble}
\pagestyle{fancy}

\begin{document}

\fancyhead[R]{\url{https://ianchen3.github.io} \\ \url{linkedin.com/in/ian-chen-b92491324/}}
\fancyhead[C]{\Huge \textbf{Ian Chen}}
\fancyhead[L]{ianchen3@illinois.edu \\ (650) 505-6879}

\section*{Education}
\vspace*{-2.5mm}

\textbf{University of Illinois} Urbana-Champaign, College of Liberal Arts and Sciences \hfill \daterange{8}{2023}{5}{2027} \\
\emph{Major in Statistics \& Computer Science} --- \emph{Cumulative GPA}: 3.96 / 4.00 \\

\vspace*{-5mm}
\rule{\textwidth}{0.1pt}
\vspace*{-10mm}

\section*{Publications}
\vspace*{-2.5mm}

\begin{itemize}
    \cramped
    \item Park, M., Yi, H., \textbf{Chen, I.}, Chacko, G., Warnow, T. (2026).
          Modeling Citations and Cartels.
          Preprint. \href{https://doi.org/10.21203/rs.3.rs-9270520/v1}{10.21203/rs.3.rs-9270520/v1}
    \item Vu-Le, T., Park, M., \textbf{Chen, I.}, Chacko, G., Warnow, T. (2025).
          Using Stochastic Block Models for Community Detection: The issue of edge-connectivity.
          Applied Network Science. \href{https://doi.org/10.1007/s41109-025-00747-2}{10.1007/s41109-025-00747-2}
    \item Vu-Le, T., Lamy, J., \textbf{Chen, I.}, Park, M., Harb, E., Chacko, G., Warnow, T. (2025).
          Dense Subgraph Clustering and a New Cluster Ensemble Method.
          Complex Networks 2025. \href{https://doi.org/10.1007/978-3-032-16719-4\_3}{10.1007/978-3-032-16719-4\_3}
\end{itemize}

\rule{\textwidth}{0.1pt}
\vspace*{-10mm}

\section*{Technical Skills}
\vspace*{-2.5mm}

\textbf{Programming Languages; Tools} \\
Python, C++, Go, C\#, R, Rust, Bash, SQL, MIPS; Git, \LaTeX \\

\vspace*{-5mm}
\rule{\textwidth}{0.1pt}
\vspace*{-10mm}

\section*{Work Experience}
\vspace*{-2.5mm}

\textbf{Siebel School of Computing and Data Science} \\
\emph{Undergraduate Research Assistant} \hfill \daterangepres{5}{2025}
\begin{itemize}
    \cramped
    \item Implemented massively parallizable community search algorithms for high performance computing systems.
    \item Conducted a literature review comparing state-of-the-art (SotA) community detection (CD) algorithms
    \item Designed experiments demonstrating that a novel CD algorithm outperforms the SotA alternatives.
    \item Edited manuscripts for two publications accepted in peer reviewed conferences and journals.
\end{itemize}

\textbf{Maven Optronics} \hfill \daterange{5}{2024}{8}{2024} \\
\emph{Software Engineer Intern}
\begin{itemize}
    \cramped
    \item Implemented Manufacturing Execution System (MES) materials traceability modules — populated mySQL tables to link raw material barcode IDs (with photo records) to assembled MiniLED boards.
    \item Developed automatic deployment of software releases via win32 installer (.msi) through administrative installation across NAS server.
\end{itemize}

\rule{\textwidth}{0.1pt}
\vspace*{-10mm}

\section*{Projects}
\vspace*{-2.5mm}

\textbf{ClassFS: course explorer filesystem} (\url{https://github.com/ianchen3/ClassFS}) \hfill \daterange{8}{2024}{12}{2024}
\begin{itemize}
    \cramped
    \item Collaborated with two team members to build virtual file system through FUSE, allowing users to browse the UIUC course catalog with their file explorer of choice.
    \item Implemented system calls hooks for \verb+open+, \verb+read+, \verb+readdir+, and \verb+stat+.
\end{itemize}

\textbf{Directle: NLP-enabled linguistics game} (\url{https://github.com/ianchen3/Directle}) \hfill \daterange{1}{2024}{5}{2024}
\begin{itemize}
    \cramped
    \item Collaborated with three team members to build a full-stack website hosting an English word game. The game guides players to guess a secret word through one-word semantic clues.
    \item Implemented a word2vec NLP model using the top 10M wikipedia articles to calculate of word similarities.
\end{itemize}

\rule{\textwidth}{0.1pt}
\vspace*{-10mm}

\section*{Relevant Coursework}
\vspace*{-2.5mm}

\textbf{Computer Science:}
\emph{Distributed Systems};
\emph{Computer Vision};
\emph{Machine Learning};
\emph{Communication Networks};
\emph{Operating Systems};
\emph{Computer Architecture I \& II};
\emph{Algorithms I \& II};
\emph{Randomized Algorithms};
\emph{Computational Biology}

\textbf{Statistics \& Math:}
\emph{Probability I \& II};
\emph{Numerical Analysis};
\emph{Graph Theory};
\emph{Differential Equations};
\emph{Real Analysis}

\textbf{Physics:}
\emph{Classical Mechanics};
\emph{E\&M};
\emph{Thermodynamics};
\emph{Quantum Mechanics};
\emph{Special Relativity}

\end{document}
